\documentclass[12pt,a4paper]{report}

% Preamble generated by build_preamble for report

% Source: sources/preamble/standard_packages.tex

% Encoding and language
\usepackage[T1]{fontenc}
\usepackage[utf8]{inputenc}
\usepackage[british]{babel}

% Mathematics
\usepackage{mathtools}
\usepackage{amsthm}
\usepackage{libertine}
\usepackage[libertine]{newtxmath}
\usepackage[mathscr]{euscript}

% General packages
\usepackage{enumitem}
\usepackage{comment}

% Figures
\usepackage{graphicx}
\usepackage{float}

% Diagrams
\usepackage{tikz-cd}
\usetikzlibrary{decorations.markings}

% Bibliography
\usepackage[
    backend=biber,
    style=alphabetic,
    maxnames=10,
    maxalphanames=10]{biblatex}
\addbibresource{refs.bib}

% Colors
\usepackage{xcolor}
\definecolor{CustomBlue}{RGB}{51,51,153}

% References and links
\usepackage{hyperref}
\hypersetup{
    colorlinks=true,
    linkcolor=CustomBlue,
    citecolor=CustomBlue,
    urlcolor=CustomBlue
}
\urlstyle{same}
\usepackage[noabbrev]{cleveref}

% Source: sources/preamble/standard_macros.tex

% Theorem-like environments
\theoremstyle{plain}
\newtheorem{theorem}{Theorem}[section]
\newtheorem*{theorem*}{Theorem}
\newtheorem{lemma}[theorem]{Lemma}
\newtheorem{proposition}[theorem]{Proposition}
\newtheorem{corollary}[theorem]{Corollary}
\newtheorem{fact}[theorem]{Fact}

\theoremstyle{definition}
\newtheorem{definition}[theorem]{Definition}
\newtheorem{construction}[theorem]{Construction}
\newtheorem{conjecture}[theorem]{Conjecture}
\newtheorem{exercise}[theorem]{Exercise}

\theoremstyle{remark}
\newtheorem{remark}[theorem]{Remark}
\newtheorem{notation}[theorem]{Notation}
\newtheorem{example}[theorem]{Example}
\newtheorem{question}[theorem]{Question}

% Alphabetic theorem items with ordinary enumerate indentation
% Printed labels: (a), (b), ...; references: 1.2.3(a), 1.2.3(b), ...
\newlist{theoremitems}{enumerate}{1}
\setlist[theoremitems,1]{
    label=(\alph*),
    ref=\thetheorem(\alph*)
}

% Equation numbering
\numberwithin{equation}{section}

% Cross-reference names
\crefname{theorem}{theorem}{theorems}
\Crefname{theorem}{Theorem}{Theorems}
\crefname{lemma}{lemma}{lemmas}
\Crefname{lemma}{Lemma}{Lemmas}
\crefname{proposition}{proposition}{propositions}
\Crefname{proposition}{Proposition}{Propositions}
\crefname{corollary}{corollary}{corollaries}
\Crefname{corollary}{Corollary}{Corollaries}
\crefname{fact}{fact}{facts}
\Crefname{fact}{Fact}{Facts}
\crefname{definition}{definition}{definitions}
\Crefname{definition}{Definition}{Definitions}
\crefname{construction}{construction}{constructions}
\Crefname{construction}{Construction}{Constructions}
\crefname{conjecture}{conjecture}{conjectures}
\Crefname{conjecture}{Conjecture}{Conjectures}
\crefname{exercise}{exercise}{exercises}
\Crefname{exercise}{Exercise}{Exercises}
\crefname{remark}{remark}{remarks}
\Crefname{remark}{Remark}{Remarks}
\crefname{notation}{notation}{notations}
\Crefname{notation}{Notation}{Notations}
\crefname{example}{example}{examples}
\Crefname{example}{Example}{Examples}
\crefname{question}{question}{questions}
\Crefname{question}{Question}{Questions}
\crefname{theoremitemsi}{item}{items}
\Crefname{theoremitemsi}{Item}{Items}

% Source: sources/preamble/todo.tex

% Warning symbol
\usepackage{utfsym}

% To-do notes
\newif\ifshowtodos
\showtodosfalse

\newcommand{\todo}[1]{%
    \ifshowtodos
        \noindent{\color{red}\usym{26A0} #1}%
    \fi
}

% Source: sources/preamble/author_comments.tex

% Author comments
\newif\ifshowcomments
\showcommentsfalse

\newcommand{\newauthorcomment}[3]{%
    \expandafter\newcommand\csname #1\endcsname[1]{%
        \ifshowcomments
            \noindent{\textcolor{#3}{\textbf{#2:} ##1}}%
        \fi
    }%
}

% Source: config/user_example.sh

\title{Example report document}
\author{Example Author\thanks{Example University\\Department of Mathematics, 1 Example Road\\Example City, Country; \href{mailto:author@example.com}{author@example.com}; \href{https://example.com}{https://example.com}}}
\date{\today}

% Source: sources/examples/report.tex

\begin{document}

\showtodostrue
\showcommentstrue

\newauthorcomment{alice}{Alice}{red}
\newauthorcomment{bob}{Bob}{red}

\maketitle
\tableofcontents

\chapter{A Sample Chapter}

\section{Structured statements}

\begin{definition}
    \label{definition:longform-object}
    A \emph{longform object} is an object \(X\) equipped with a map
    \[
        \alpha_X \colon X \to G(X)
    \]
    satisfying the following properties.
    \begin{theoremitems}
    \item The object \(X\) exists in the category under consideration.
        \label{item:longform-existence}
    \item It is unique up to unique isomorphism.
        \label{item:longform-uniqueness} 
    \end{theoremitems}
\end{definition}

\begin{proposition}
    \label{proposition:longform-main}
    Every longform object is determined by the data in
    \Cref{item:longform-existence,item:longform-uniqueness}.
\end{proposition}

\begin{proof}
    This is a template proof. The point is that the item references are attached to
    the theorem-like environment \Cref{definition:longform-object}, so their labels
    are numbered using the surrounding theorem number.
\end{proof}

\begin{remark}
    \label{remark:longform-citations}
    Long-form documents often contain many references.
    This example cites Hartshorne \cite{har77} and the Stacks Project \cite{stacks}.
\end{remark}

\section{Diagrams and ordinary lists}

The template supports commutative diagrams. For instance, one may write
\[
    \begin{tikzcd}
        A \arrow[r, "\phi"] \arrow[d, "\alpha"'] &
        B \arrow[d, "\beta"] \\
        C \arrow[r, "\psi"'] &
        D .
    \end{tikzcd}
\]
Such diagrams are useful in lecture notes and book-length projects.

The template also supports ordinary enumerated lists:
\begin{enumerate}[label=(\alph*), ref=(\alph*)]
    \item A first ordinary item.
        \label{item:longform-ordinary-first} 
    \item A second ordinary item.
        \label{item:longform-ordinary-second} 
\end{enumerate}
We can refer to \Cref{item:longform-ordinary-second} as an ordinary item, and
to \Cref{item:longform-uniqueness} as a theorem-numbered item.

\chapter{Another Chapter}

\section{Resetting theorem numbers}

This chapter checks that theorem-like numbering behaves correctly after a
chapter change.
\todo{Here is a to-do note.}

\begin{lemma}
    \label{lemma:second-chapter}
    The references from the first chapter, such as
    \Cref{definition:longform-object}, still work from later chapters.
\end{lemma}

\begin{proof}
    This proof is only illustrative. It also demonstrates a displayed equation:
    \begin{equation}
        \label{equation:pull-back}
        (f \circ g)^* = g^* \circ f^* .
    \end{equation}
\end{proof}

\alice{We should add an example here.}

\begin{example}
    \label{example:external-link}
    A long-form document may also contain external links, for instance to
    \href{https://www.ctan.org/pkg/biblatex}{the biblatex package page}.
\end{example}

\bob{I've added the example above.}

\printbibliography

\end{document}
